% Generated by openai gpt-4.1 on 2026-11-09T13:50:34.032935+00:00, prompt geosmile-hull-v9 + card-mode addendum, temperature 0 (generate_card_narrative.py --stakeholder planner). Checks: passed. Do not edit by hand.
\renewcommand{\cardHeadline}{The model relies most on \kt{Old Town}, but built-environment features matter most in \kt{Longhill} and \kt{Bilton Grange}.}
\renewcommand{\cardSecA}{The model is most sensitive to \kt{Old Town} (029E, rank \pv{1/168}), ahead of Dairycoates and Sculcoates. In Old Town, it relies overwhelmingly on \kt{recent crime history} (\pv{82\%} of local importance); built-environment features contribute only \pv{9\%}, mainly \kt{nightlife}.}
\renewcommand{\cardSecB}{Across Hull, \kt{crime history} dominates model attribution (\pv{64\%}), followed by census characteristics (\pv{18\%}); built-environment features account for \pv{14\%}, led by \kt{green space} and roads. Built-environment influence peaks in \kt{Longhill} (\pv{50\%}, green space) and \kt{Bilton Grange} (\pv{50\%}, transit stops).}
\renewcommand{\cardFooter}{Attribution shares reflect model behaviour, not planning targets.}
\renewcommand{\cardModel}{gpt-4.1}
\renewcommand{\cardTileA}{\tile{1/168}{GeoSensitivity rank, Old Town}}
\renewcommand{\cardTileB}{\tile{82\%}{Crime history, Old Town}}
\renewcommand{\cardTileC}{\tile{14\%}{Built environment, Hull-wide}}
\renewcommand{\cardTileD}{\tile{50\%}{Highest built share, Longhill}}
